%% chemfrag-polymers.code.tex -- polymers (repeat units in brackets)
%%
%% Copyright (C) 2026 Vaibhav Blayer
%% LaTeX Project Public License 1.3c or later (see chemfrag.sty).
%%
%% \chemfragpolymer[<keys>]{<repeat unit>} draws
%%     <left> -[@{op}] <repeat unit> -[@{cl}] <right>
%% with \chemfig and puts chemfig's \polymerdelim brackets on the two bonds
%% that cross them.  <repeat unit> is either the name of a unit defined with
%% \chemfragdefinepolymer (polyethylene, PVC, nylon66, ...) or chemfig code.
%% Like \chemmove, the brackets need two compilations.

\keys_define:nn { chemfrag / polymer }
  {
    n           .tl_set:N = \l__chemfrag_poly_n_tl ,
    n           .initial:n = n ,
    height      .tl_set:N = \l__chemfrag_poly_height_tl ,
    height      .initial:n = 6pt ,
    depth       .tl_set:N = \l__chemfrag_poly_depth_tl ,
    depth       .initial:n = 6pt ,
    delimiters  .tl_set:N = \l__chemfrag_poly_delims_tl ,
    delimiters  .initial:n = {[]} ,
    left        .tl_set:N = \l__chemfrag_poly_left_tl ,
    right       .tl_set:N = \l__chemfrag_poly_right_tl ,
    bond        .tl_set:N = \l__chemfrag_poly_bond_tl ,
    bond        .initial:n = 1 ,
    chemfig     .tl_set:N = \l__chemfrag_poly_chemfig_tl ,
  }

\int_new:N \g__chemfrag_poly_int

%% \chemfragdefinepolymer{<name>}[<default keys>]{<repeat unit>}
\NewDocumentCommand \chemfragdefinepolymer { m O{} m }
  {
    \tl_gclear_new:c { g__chemfrag_poly_#1_unit_tl }
    \tl_gset:cn { g__chemfrag_poly_#1_unit_tl } {#3}
    \tl_gclear_new:c { g__chemfrag_poly_#1_keys_tl }
    \tl_gset:cn { g__chemfrag_poly_#1_keys_tl } {#2}
  }

\NewDocumentCommand \chemfragpolymer { O{} m }
  {
    \group_begin:
      \int_gincr:N \g__chemfrag_poly_int
      \tl_set:Ne \l__chemfrag_poly_name_tl { \tl_to_str:n {#2} }
      \cs_if_exist:cTF { g__chemfrag_poly_ \l__chemfrag_poly_name_tl _unit_tl }
        {
          \keys_set:nv { chemfrag / polymer }
            { g__chemfrag_poly_ \l__chemfrag_poly_name_tl _keys_tl }
          \keys_set:nn { chemfrag / polymer } {#1}
          \__chemfrag_polymer:v
            { g__chemfrag_poly_ \l__chemfrag_poly_name_tl _unit_tl }
        }
        {
          \keys_set:nn { chemfrag / polymer } {#1}
          \__chemfrag_polymer:n {#2}
        }
    \group_end:
  }
\tl_new:N \l__chemfrag_poly_name_tl
\cs_generate_variant:Nn \keys_set:nn { nv }

%% The open/close nodes are points on the two crossing bonds (a quarter of
%% the bond from its outer end, leaving a short stub outside), named
%% cfpolyo<k>/cfpolyc<k> with a counter so that several polymers can share a
%% paragraph.
\cs_new_protected:Npn \__chemfrag_polymer:n #1
  {
    \use:e
      {
        \exp_not:N \chemfig
          [ \exp_not:V \l__chemfrag_poly_chemfig_tl ]
          {
            \tl_if_empty:NTF \l__chemfrag_poly_left_tl
              { \exp_not:N \vphantom { \exp_not:n { CH } } }
              { \exp_not:V \l__chemfrag_poly_left_tl }
            - [ \c_atsign_str { cfpolyo \int_use:N \g__chemfrag_poly_int , .25 }
                , \l__chemfrag_poly_bond_tl ]
            \exp_not:n {#1}
            - [ \c_atsign_str { cfpolyc \int_use:N \g__chemfrag_poly_int , .75 }
                , \l__chemfrag_poly_bond_tl ]
            \exp_not:V \l__chemfrag_poly_right_tl
          }
        \exp_not:N \polymerdelim
          [
            delimiters = { \exp_not:V \l__chemfrag_poly_delims_tl } ,
            height = \l__chemfrag_poly_height_tl ,
            depth = \l__chemfrag_poly_depth_tl ,
            indice = { \exp_not:V \l__chemfrag_poly_n_tl }
          ]
          { cfpolyo \int_use:N \g__chemfrag_poly_int }
          { cfpolyc \int_use:N \g__chemfrag_poly_int }
      }
  }
\cs_generate_variant:Nn \__chemfrag_polymer:n { v }
